%% ijcai ecai 2025
%% 7 pages + 2 pages for ethics, ack, refs

\documentclass{article}
\pdfpagewidth=8.5in
\pdfpageheight=11in

% The file ijcai25.sty is a copy from ijcai22.sty
% The file ijcai22.sty is NOT the same as previous years'
\usepackage{ijcai25}

% Use the postscript times font!
\usepackage{times}
\usepackage{soul}
\usepackage{url}
\usepackage[hidelinks]{hyperref}
\usepackage[utf8]{inputenc}
\usepackage[small]{caption}
\usepackage{graphicx}
\usepackage{amsmath}
\usepackage{amsthm}
\usepackage{booktabs}
\usepackage{algorithm}
\usepackage{algorithmic}
\usepackage[switch]{lineno}

% Comment out this line in the camera-ready submission
\linenumbers

\urlstyle{same}

% the following package is optional:
%\usepackage{latexsym}

% See https://www.overleaf.com/learn/latex/theorems_and_proofs
% for a nice explanation of how to define new theorems, but keep
% in mind that the amsthm package is already included in this
% template and that you must *not* alter the styling.
\newtheorem{example}{Example}
\newtheorem{theorem}{Theorem}

% Following comment is from ijcai97-submit.tex:
% The preparation of these files was supported by Schlumberger Palo Alto
% Research, AT\&T Bell Laboratories, and Morgan Kaufmann Publishers.
% Shirley Jowell, of Morgan Kaufmann Publishers, and Peter F.
% Patel-Schneider, of AT\&T Bell Laboratories collaborated on their
% preparation.

% These instructions can be modified and used in other conferences as long
% as credit to the authors and supporting agencies is retained, this notice
% is not changed, and further modification or reuse is not restricted.
% Neither Shirley Jowell nor Peter F. Patel-Schneider can be listed as
% contacts for providing assistance without their prior permission.

% To use for other conferences, change references to files and the
% conference appropriate and use other authors, contacts, publishers, and
% organizations.
% Also change the deadline and address for returning papers and the length and
% page charge instructions.
% Put where the files are available in the appropriate places.


% PDF Info Is REQUIRED.

% Please leave this \pdfinfo block untouched both for the submission and
% Camera Ready Copy. Do not include Title and Author information in the pdfinfo section
\pdfinfo{
/TemplateVersion (IJCAI.2025.0)
}

\title{Selecting Images}

\iffalse
\author{
Katerina Oikonomou \and Stasinos Konstantopoulos
\affiliations
Institute of Informatics and Telecommunications
\\ NCSR `Demokritos',
Ag. Paraskevi, Greece
\emails
\{first, second\}@iit.demokritos.gr
}
\else
\author{
Anonymous
\affiliations
Affiliation
\emails
anonymous@affiliation
}
\fi

\begin{document}

\maketitle

\begin{abstract}


\end{abstract}

\section{Introduction}

% Positioning. Background.
Sustainable AI has emerged in an era where AI workload enormously increase the data center emissions~\cite{jha-etal:2025}. 
Computationally expensive and time-consuming neural networks are being trained daily with massive datasets, while the demand for even larger, better or more task-oriented datasets shows no sign of slowing down.
Questiong like: "How much more data do I need~\cite{mahmood-etal:2022}?" and "do we need more training data~\cite{zhu-vondrick-etal:2016}?" highlight that large-scale datasets have increasingly become a contraint rathen than a beneficial resource. 

% Here about DD
Towards that, dataset distillation (DD) has been proposed as a remedy~\cite{lei-tao:2024}. 
DD aprroaches synthesize smaller synthetic datasets that preserve the important representation characteristics of their original counterpart and at the same time achieve similar results. 
However these methods still generate new data in an ecosystem already burdened by data overload and significant computational cost.


In this work, we take a different stance and investigate whether sustainable AI can be advanced through coreset selection rather than dataset synthesis. 
We argue that it is feasible to capture specific instances, within a training dataset, that are the most informative ones and preserve the networks decision information, resulting faster and thus more energy efficient training.
We analyze the trade-off between subset size, accuracy, and training cost and show that carefully selected coreset can rival large-scale datasets in termes of accuracy. 
 


\section{Method}

The main idea of  the proposed approach is that in a contrastive learning approach we want to maximize the value we 
get from the labelling. 

We hypothesize that the representative instances i.e the ones closest to the classes centroid are not informed 
of the classes true extend and even more so of the nuances of the boundaries between this class and the other classes.

Based on the above we claim that coreset selection should be biased towards the instances further away from the centroid, without completely discarding representative instances. 




To establish the validity of this hypothesis we developed the following method:



\subsection{Variational Autoencoder}


The autoencoders are used for non-linead dimentionality reduction while preserving the important characteristics of each sample.
Variational autoencoders provide a propabilistic latent representation. 

Here we utilized the pre-trained SD VAE which is trained on the COCO dataset and encodes all the important information into an informative comporessed latent space.

Instead of performing instance selection directly at the raw VAE's latent space we trained an embedding network on top of the VAE with a discriminative head to produce features.

\subsection{Feature Discrimination}

Feature discrimination is an important challenge in computer vision..
It is important that embeddings exhibit high discrimination along different classes as well as intra-class compactness. 
Towards that the authors at~\cite{kansizoglou-santavas-etal:2020} proposed a class hyperplane-based segregation on the contrary of the common class centers separation scheme.


\subsection{Coreset selection}

At~\cite{doshi-rinal-bhadka-etal:2013} the authors proposed a data clustering mining algorithm i.e a fast and greedy algorithm to select k points from a cluster. 
The first center is selected randomly the second is greedily select the point furthest from the first.
Each subsequent center is selected as the point that maximize the minimum distance to the closest already selected centers. 
Thus this algorithm has a bias to the boundaries of the cluster and then choose instances from the inside in order to evenly choose intances for the cluster.

Furthermore, at~\cite{khan-kafi-sarker-etal:2021} the authors utilize cosine similarity to assign data to clusters showing its benefits compared with the euclidean one.
Towards that we adopted the same cosine metric to the above methodolody of the furthers-first clustering approach.


\section{Embedding approach}

In this work we propose an embedding pipeline that combines a frozen variational autoencoder with a lightweight discriminative network and a hyperplane-based classification head.
The goal is to obtain a feature space that holds all the necessary encoded information from the VAE, while at the same time instances from the same class form tight clusters, while different classes are well separated.
This space form the basis for instance selection.


Given an input imageInitially 


Use synthetic data to demonstrate how it works.


\section{Experiments and Results}

Related work.

\begin{table}[h!]
    \centering
    \caption{Test Accuracy at Various Data Percentages on Resnet-18.}
    \label{tab:performance_data}
    \begin{tabular}{lccccccccc}
        \toprule
        \textbf{Approach} & \textbf{10\%} & \textbf{20\%} & \textbf{30\%} & \textbf{40\%} & \textbf{50\%} & \textbf{60\%} & \textbf{70\%} & \textbf{80\%} & \textbf{90\%} \\
        \midrule
        GradMatch~\cite{killamsetty-ksrishnateja-etal:2021} &  \\
        Ours &  \\
        \bottomrule
    \end{tabular}
\end{table}

\begin{table}[h!]
    \centering
    \caption{Test Accuracy at Model-Agnostic Coreset Selection on CIFAR10.}
    \label{tab:model_agnostic}
    \begin{tabular}{lccccccccc}
        \toprule
        \textbf{Model} & \textbf{10\%} & \textbf{20\%} & \textbf{30\%} & \textbf{40\%} & \textbf{50\%} & \textbf{60\%} & \textbf{70\%} & \textbf{80\%} & \textbf{90\%} \\
        \midrule
        Custom Convnet & 53.09 & 61.81 & 65.71 & 68.65 & 70.07 & 72.24 & 73.79 & 74.99 & 76.32 \\
        Resnet-18 & 79.21& 86.22 & 89.59 & 91.19& 92.42& 93.19 & 93.56& 93.58& 94.38 \\
        Resnet-50 & 72.32 & 84.87 & 88.21 & 90.98& 91.63 & 92.54 & 93.06 & 93.65&94.10 \\
        Resnet-101& 74.47 & 85.92 & 89.59 & 91.37& 92.66 & 93.62 & 93.98& 94.44& 94.79 \\
        Densenet-169 &  76.74 & 86.48 & 89.65 &91.62 &92.69& 93.22& 94.34&94.72 & 94.93 \\

        \bottomrule
    \end{tabular}
\end{table}

\begin{table}[h!]
    \centering
    \caption{Training time at Model-Agnostic Coreset Selection on CIFAR10 in minutes.}
    \label{tab:training_time}
    \begin{tabular}{lccccccccc}
        \toprule
        \textbf{Model} & \textbf{10\%} & \textbf{20\%} & \textbf{30\%} & \textbf{40\%} & \textbf{50\%} & \textbf{60\%} & \textbf{70\%} & \textbf{80\%} & \textbf{90\%} \\
        \midrule
        Custom Convnet &  15.16 & 20.42 & 25.99 & 31.48 & 36.96 & 42.27 &47.77 & 53.24 & 58.87 \\
        Resnet-18 &  23.70 & 30.27& 37.30 & 44.10& 51.16& 58.01& 65.02& 42.71& 80.23\\
        Resnet-50 &  14.67& 20.80&  26.75& 32.82& 38.82 & 44.83& 50.67 & 72.81& 112.16\\
        Resnet-101 & 37.29  &56.98  & 76.58 & 95.81& 114.91& 134.52& 153.47 & 173.08& 192.65 \\
        Densenet-169 & 37.39 & 57.86& 78.43& 98.43& 118.87& 138.99& 159.34& 180.04& 200.33\\

        \bottomrule
    \end{tabular}
\end{table}
Experimental setup.

% https://github.com/Spijkervet/SimCLR/

Results and discussion.

\section{Discussion}

The observed that the choosen boundaries reflect the dimensionality of the dataset.

\section{Conclusions}

Wow!


\appendix

\section*{Ethical Statement}

There are no ethical issues.

\iffalse
\section*{Acknowledgments}

MANOLO
PharosAI

\fi

%% The file named.bst is a bibliography style file for BibTeX 0.99c
\bibliographystyle{named}
\bibliography{ijcai25}

\end{document}

